\section{Additional arguments, implementation scope, and retained records}
\label{sec:scope}
\subsection{Why a whole-graph zero-degree fallback fails}
\begin{proposition}[A whole-graph zero-degree fallback is invalid in the
sparse regime]\label{restore:fallback}
Consider a balanced two-group SBM with within probability $p$, between
probability $p/2$, and $d=np=\log\log n$. The probability of at least
one isolated vertex tends to one. For every deterministic fallback
partition fixed in advance, there is a balanced true partition for which
an algorithm returning that fallback whenever an isolated vertex exists
has expected permutation-invariant misclassification bounded below by
$1/4+o(1)$. Its local oracle risk tends to zero.
\end{proposition}

\begin{proof}
For every vertex the isolation probability is
$(1-p)^{n/2-1}(1-p/2)^{n/2}
=\exp\{-3d/4+o(1)\}$. If $N_0$ counts isolated vertices, then
$\E N_0=n\exp\{-3d/4+o(1)\}\to\infty$. For distinct vertices,
their two isolation events share only the edge between them, giving
\[
 \frac{\operatorname{Var}(N_0)}{(\E N_0)^2}
 \le \frac1{\E N_0}+\frac p{1-p}\longrightarrow0.
\]
Choose a balanced true partition whose overlap with the fixed fallback
has error at least $1/4$ after either label permutation. Such a partition
exists by splitting each fallback class approximately evenly. The loss
lower bound follows. Standard incident-row Chernoff comparison gives
oracle risk at most $Ce^{-cd}\to0$ on this separated model.
\end{proof}

Thus the logarithmic restriction cannot be removed merely by deleting
text from the unregularized theorem. The valid repair is to restore a
positive floor and local risk accounting. With the regularized encoder,
a zero-degree vertex has a zero reconstructed row and causes no
whole-graph fallback.

\subsection{The existing degree-plus-spectrum action remains valid}

Write $\mathbf d=A\one$ for the raw degree vector. For any strictly
positive diagonal matrix $W$, retain
\[
 WAWU=U\Lambda,\qquad U^\top U=I,\qquad
 X=[\one,WU],\qquad Y=[\mathbf d,W^{-1}U\Lambda].
\]
Then the exact identities
\begin{equation}\label{eq:sparse-old-action}
 AX=Y,\qquad YX^\dagger=A\Pi_X,\qquad A\Pi_X\one=\mathbf d
\end{equation}
hold for every graph, including graphs with zero degrees. Here $\Pi_X$
is the Euclidean orthogonal projector onto the column space of $X$;
no full-column-rank assumption is necessary. The earlier choice
$W_{ii}=\min(1,2\bar D/D_i)$ for $D_i>0$, with $W_{ii}=1$ for
$D_i=0$ and $\bar D=n^{-1}\sum_iD_i$, supplies a positive reweighting
when $\bar D>0$. The all-zero graph can be assigned a fixed output.
The block-count error is exactly $-A(I-\Pi_X)Z$.

These facts need no logarithmic expected-degree assumption. They are
computational identities, however: their validity does not give the
sharp tail of the remaining random block-count error. In DC-SBM a
normalization must also preserve the effective likelihood coefficients,
as opposed to preserving only a population mean.


\subsection{Algorithm distinctions and scope}
The current ordinary article uses the regularized block Bernoulli LRT,
with fixed degree parameters equal to one and no target profiling.
The current DC result in this companion uses clamped, group-normalized
nuisance estimates and profiles the target separately for each candidate.
Both retain raw degrees and the informative normalized eigensystem,
apply exact degree correction and simple hollowing, and control finite
hard refits through a uniform invariant basin.

The older ordinary free-profile growing/replacement/soft-EM schedule
in \cref{sec:sbm-proof} is a different finite algorithm. Its separate
proof also covers probabilities outside bounded $np^2$. The old common-box
proof uses a simultaneous entry bound and must not be attached without
change to the present sparse encoder. In particular, exact degree
preservation together with $|R_{ij}|\le Cp$ would force all raw degrees
to be at most $Cd$, which is not a valid sparse uniform premise.
The logarithmic restrictions written in the historical part describe
those earlier unregularized proofs, not the current theorem.

The earlier DC reconstruction $AH_D$ is different from the Euclidean
$A\Pi_X$, even though both can be evaluated from the same retained
input. Earlier experimental curves do not automatically evaluate either
the new empirical floor or its clamped nuisance fit. The present claims
are mathematical risk bounds, not conclusions inferred from simulations.
Numerical eigensolver tolerances and optimization tolerances require their
own control of induced score errors; exact arithmetic defines these
algorithms. The largest absolute eigenvalues are retained so that a
full-rank indefinite signal is included.

The constants may depend on the fixed signal gap, group proportions,
bounded degree-parameter range, and $C_1$. A bounded risk ratio is not
a claim that the ratio tends to one. General DC-SBM with unbounded
$np^2$, a vanishing signal gap, or degree parameters approaching zero
is outside the present theorem. The general known-target-degree
benchmark is obstructed by \cref{thm:imposs}; equal exposures retain
the positive result.

\subsection{Sparse proof history and review records}
The 9 October degree-plus-spectrum draft contained the sparse truncated
row-energy mechanism at $d\to\infty$, $d=o(\log n)$. Its score errors
were $O(\sqrt d)$ or $o(d)$ and therefore established a leading exponent,
not the present constant risk ratio. The current argument keeps its
same-threshold mechanism, upgrades the local full-to-deleted score
comparison and fitted contrast errors to $O(1)$, and preserves the sharp
$d^{-1/2}$ prefactor. The earlier source remains in the repository at
\href{https://github.com/zhixin0825/spectral-lrt/blob/de60b64c7d7f127ae6dc1d010c840af707a080da/paper/previous_draft/sbm_spectral_lrt_20261009.tex}{commit \texttt{de60b64}}.

The source package includes the \texttt{audits/} directory and its index.
It retains the sparse-history audit, the independent local-floor LOO
review, the same-threshold tail and conditioning review, the high-scale
probability review, the ordinary algorithm-scope review, the final
combined review, and earlier DC and regularization audit snapshots.
These are internal working checks, not external peer review. Statements
of then-open gaps in older snapshots are superseded by the current
proofs. The full pre-split revision is preserved separately in the
repository; it is an archive rather than the reading order of these notes.

An inventory in \texttt{README.md} records where each old source module
now belongs. The original implementation/audit narrative is retained as
\texttt{archive/scope\_audit.tex}; the obsolete end-of-revision status
narrative is retained as \texttt{archive/sparse\_interfaces.tex}. These
source records are not assertions of a remaining logarithmic lower bound.
